\section{A scalar entropy inequality}
\label{sec:entropy}

The stability argument needs a scalar criterion that forces an
empirical field law toward a Gaussian law.  We first bound relative
entropy below by a squared discrepancy involving the field and its
hyperbolic tangent.  We then identify the Gaussian law compatible
with the scalar consistency equations.  Together with the equality
case of the first estimate, this characterizes the consistent
equality laws.

Throughout this section, $N(d,s)$ denotes the Gaussian law with mean
$d$ and \emph{variance} $s$.  All scalar expectations are taken with
respect to the indicated probability law.

\begin{lemma}[Scalar entropy inequality]
\label{lem:scalar-entropy}
Let $0<j<1$, $d\in\R$, and $s>0$.  Let $P$ be a probability law on $\R$
with finite second moment, and set
\[
 m(y)=\tanh y,\qquad q=\E_P m^2,\qquad b=1-q,\qquad
 a=\E_P[(y-d)m],\qquad S=s+jq.
\]
If $s\ge jq$, then
\begin{equation}
 D(P\Vert N(d,s))\ge \frac{j(a-sb)^2}{2sS}.
 \label{eq:scalar-entropy}
\end{equation}
Equality holds if and only if $P=N(d,s)$.
\end{lemma}

\begin{proof}
The second-moment assumption makes the right-hand side finite, so
it suffices to consider the case $D:=D(P\Vert N(d,s))<\infty$.  In this case $P$ is
absolutely continuous, and hence $q>0$: otherwise $P$ would be the
point mass at zero.  Also $q<1$, because $\tanh^2 y<1$ for every finite
$y$.  We use the following dimensionless parameters:
\begin{equation}
 k=\frac{jq}{s}\in(0,1],\qquad C=\frac{ja}{s},\qquad
 B=jb\in(0,1),\qquad a_0=1-B\in(0,1).
 \label{eq:entropy-scalars}
\end{equation}
In these parameters, multiplying the desired inequality by $j$
gives the target
\begin{equation}
 jD\ge \frac{(C-B)^2}{2(1+k)}.
 \label{eq:entropy-target}
\end{equation}

We will combine two entropy bounds.  A monotone change of variables
handles one range of the discrepancy $C-B$; comparison with a Gaussian
of a different variance handles the remaining range.  We retain
strictness in both arguments to identify all equality cases.

\paragraph{A change-of-variables bound.}
For the change of variables, fix $u<1$.  The map
\[
 T_u(y)=y-u\tanh y
\]
is an increasing $C^1$ diffeomorphism of $\R$ onto itself, with derivative
$T_u'(y)=1-u(1-m(y)^2)>0$.  Change of variables in relative entropy yields
\begin{align*}
 D
 &=D((T_u)_\#P\Vert N(d,s))
   +\frac{2ua-u^2q}{2s}
   +\E_P\log\bigl(1-u(1-m^2)\bigr)\\
 &\ge \frac{2ua-u^2q}{2s}
   +\E_P\log\bigl(1-u(1-m^2)\bigr).
\end{align*}
The entropy identity is well defined: for fixed $u<1$, $\log T_u'$
is bounded, and the second-moment assumption makes the difference
of the Gaussian log densities integrable.  To estimate the logarithmic
term, apply concavity on the segment joining $1$ and $1-u$.  For
$0\le v\le1$, this gives
\[
 \log(1-uv)\ge v\log(1-u).
\]
Consequently, with $\ell(u)=-\log(1-u)-u$,
\begin{equation}
 jD\ge u(C-B)-\frac{k u^2}{2}-B\ell(u),\qquad u<1.
 \label{eq:entropy-transport}
\end{equation}
Completing the square isolates the amount by which this bound can
exceed the target in \eqref{eq:entropy-target}:
\begin{equation}
 jD-\frac{(C-B)^2}{2(1+k)}
 \ge h(u)-\frac{\bigl(C-B-(1+k)u\bigr)^2}{2(1+k)},
 \qquad h(u)=\frac{u^2}{2}-B\ell(u).
 \label{eq:entropy-slack}
\end{equation}

We first locate a range in which the extra term $h(u)$ is strictly
positive:
\begin{equation}
 h(u)>0\quad\text{for }u\le a_0,\ u\ne0,
 \qquad h(a_0)\ge\frac{a_0^3}{6}.
 \label{eq:entropy-slack-positive}
\end{equation}
For the negative range, integrate $v/(1+v)\le v$ over
$0\le v\le -u$ to obtain $\ell(u)\le u^2/2$.  Since $B<1$, this
proves strict positivity for $u<0$.  For the positive range, use
$\ell(u)=\sum_{r\ge2}u^r/r$: it shows that $\ell(u)/u^2$ is increasing
on $0<u<1$.
Moreover, using $B=1-a_0$,
\[
 h(a_0)=\frac{a_0^2}{2}-(1-a_0)\ell(a_0)
       =\sum_{r\ge3}\frac{a_0^r}{r(r-1)}
       \ge\frac{a_0^3}{6}.
\]
It follows that $h(u)/u^2\ge h(a_0)/a_0^2>0$ for $0<u\le a_0$,
proving \eqref{eq:entropy-slack-positive}.

We now choose the transport parameter.  If
$u_*=(C-B)/(1+k)\le a_0$ and $C\ne B$, set $u=u_*$ in
\eqref{eq:entropy-slack}.  The square vanishes and $h(u_*)>0$, so
\eqref{eq:entropy-target} is strict.  If $u_*>a_0$, put
\[
 C_0=1+ka_0,
 \qquad C_1=C_0+\sqrt{\frac{1+k}{3}}\,a_0^{3/2}.
\]
In this case $C>C_0$.  The endpoint choice $u=a_0$ in
\eqref{eq:entropy-slack} still gives
\[
 jD-\frac{(C-B)^2}{2(1+k)}
 \ge \frac{a_0^3}{6}-\frac{(C-C_0)^2}{2(1+k)}.
\]
This proves strictness throughout $C_0<C<C_1$ and leaves only
the range beginning at $C_1$.

\paragraph{The remaining range.}
For $C\ge C_1$, we use the second moment to obtain the remaining
entropy bound.  Put $v=\E_P[(y-d)^2]/s>0$.  Comparison with the
Gaussian $N(d,sv)$ gives the exact identity
\[
 D=D(P\Vert N(d,sv))+\frac{v-1-\log v}{2},
\]
and hence $2D\ge v-1-\log v$.  Cauchy--Schwarz gives
\[
 v\ge\frac{a^2}{sq}=\frac{C^2}{jk}>1,
\]
where the last inequality follows from $C\ge C_1>C_0>1$ and $jk<1$.
The function $v\mapsto v-1-\log v$ is increasing for $v\ge1$, so
\[
 2jD\ge \frac{C^2}{k}-j-j\log\frac{C^2}{jk}
       \ge \frac{C^2}{k}-1-\log\frac{C^2}{k}.
\]
The second inequality removes $j$ from the lower bound.  Indeed,
with $C,k$ held fixed, the derivative of
$C^2/k-\lambda-\lambda\log(C^2/(\lambda k))$ is
$-\log(C^2/(\lambda k))\le0$ for $j\le\lambda\le1$.
Thus the remaining comparison with the target reduces to
\[
 f(C):=\frac{C^2}{k}-1-\log\frac{C^2}{k}
       -\frac{(C-B)^2}{1+k}>0
 \quad(C\ge C_1).
\]
We estimate $f$ at $C_0$ and control its increase from $C_0$ to
$C_1$.  Direct simplification at $C_0$ gives
\begin{align*}
 f(C_0)
 &=\frac1k-1+\log k
   +2a_0-a_0^2-2\log(1+ka_0)\\
 &\ge -\frac{2a_0^3}{3}.
\end{align*}
Indeed $1/k-1+\log k\ge0$ for $0<k\le1$, and
\[
 \log(1+ka_0)\le\log(1+a_0)
 \le a_0-\frac{a_0^2}{2}+\frac{a_0^3}{3}.
\]
The last bound follows by integrating
$(1+x)^{-1}\le1-x+x^2$ from zero to $a_0$.  Differentiation also gives
\begin{align*}
 f'(C_0)&=\frac2k-\frac{2}{1+ka_0}
          \ge\frac{2a_0}{1+a_0},\\
 f''(C)&=\frac{2}{k(1+k)}+\frac{2}{C^2}\ge1
          \qquad(C>0).
\end{align*}
For completeness, the first inequality follows from
\[
 \frac1k-\frac1{1+ka_0}-\frac{a_0}{1+a_0}
 =\frac{(1-k)(1+a_0+ka_0^2)}{k(1+ka_0)(1+a_0)}\ge0.
\]
The derivative bounds compensate for the possible negative value
at $C_0$.  Write $\Delta=C_1-C_0\ge a_0^{3/2}/\sqrt3$ and integrate
them to obtain
\begin{align*}
 f(C_1)
 &\ge -\frac{2a_0^3}{3}
       +\frac{2a_0}{1+a_0}\Delta+\frac{\Delta^2}{2}\\
 &\ge a_0^3\left(
       -\frac23+\frac{2}{\sqrt3(1+a_0)\sqrt{a_0}}+\frac16
       \right)>0.
\end{align*}
The strict inequality uses $0<a_0<1$ and
$1/\sqrt3+1/6>2/3$.  Since $f'$ is positive at $C_0$ and increasing
thereafter, $f(C)>0$ for every $C\ge C_1$.

The two ranges together prove strictness in
\eqref{eq:entropy-target} whenever $C\ne B$.  To finish the equality
case, suppose $C=B$.  The target then has right-hand side zero, so
equality is equivalent to $D=0$, or $P=N(d,s)$.  Conversely, Gaussian
integration by parts gives $a=s\E_P\sech^2 y=sb$ for this law,
which verifies equality.
\end{proof}

We next impose the scalar consistency equations used in the
application.  Here the Gaussian parameters themselves depend on
$P$: for given $t\ge0$ and $\sigma\ge0$, set
\begin{equation}
 M=\E_P\tanh y,\qquad q=\E_P\tanh^2 y,\qquad
 d=t+jM,\qquad s=t+\sigma^2+jq.
 \label{eq:scalar-consistency}
\end{equation}
These equations give $s\ge jq$, the comparison required by
Lemma~\ref{lem:scalar-entropy}.  When $s>0$, that lemma restricts
equality to Gaussian laws.  The next result identifies which
Gaussian laws satisfy the consistency equations when $\sigma=0$,
and also treats the degenerate case at $t=0$.

\begin{lemma}[The consistent Gaussian law]
\label{lem:gaussian-equality}
Fix $0<j<1$.  For each $t\ge0$, the equation
\begin{equation}
 s_*(t)=t+j\E_{N(s_*(t),s_*(t))}\tanh y
 \label{eq:scalar-fixed-point}
\end{equation}
has a unique nonnegative solution, where $N(0,0)$ means the point mass
at zero.  Its associated Gaussian law
$P_*(t)=N(s_*(t),s_*(t))$ is the unique Gaussian law satisfying
\eqref{eq:scalar-consistency} with $\sigma=0$.  Thus, among consistent
laws with $s>0$, equality in \eqref{eq:scalar-entropy} holds precisely at
$P_*(t)$.

The function $t\mapsto s_*(t)$ is Lipschitz with constant $(1-j)^{-1}$,
and
\begin{equation}
 t\le s_*(t)\le\frac{t}{1-j},\qquad
 q_*(t):=\E_{P_*(t)}\tanh^2 y
       =\frac{s_*(t)-t}{j}\le\frac{t}{1-j}.
 \label{eq:scalar-fixed-point-bounds}
\end{equation}
In particular $s_*(0)=q_*(0)=0$, and $s_*(t),q_*(t)>0$ for $t>0$.
\end{lemma}

\begin{proof}
We first show that a consistent Gaussian must have equal mean and
variance.  We then solve the resulting one-variable fixed-point
equation and prove the bounds in the statement.

Two Gaussian symmetry identities provide the first step.  For $s>0$,
the density $\rho_s$ of $N(s,s)$ satisfies
\[
 \frac{\rho_s(y)-\rho_s(-y)}{\rho_s(y)+\rho_s(-y)}=\tanh y.
\]
Pairing the integrals over $y$ and $-y$ therefore gives
\begin{equation}
 \E_{N(s,s)}\tanh y=\E_{N(s,s)}\tanh^2 y.
 \label{eq:gaussian-nishimori}
\end{equation}
Second, if $d\ge0$ and $s>0$, pairing the positive and negative
half-lines shows that
\begin{equation}
 \E_{N(d,s)}[\tanh y\,\sech^2 y]\ge0.
 \label{eq:gaussian-odd-sign}
\end{equation}
Here the integrand is odd and nonnegative for $y\ge0$, while the density
at $y$ is at least its value at $-y$ when $d\ge0$.

Now suppose a Gaussian law $N(d,s)$ with $s>0$ satisfies
\eqref{eq:scalar-consistency} with $\sigma=0$.  We begin by locating
its mean.  For fixed $s$, the map
\[
 d\longmapsto t+j\E_{N(d,s)}\tanh y
\]
is a contraction with Lipschitz constant at most $j$, because its
derivative is $j\E_{N(d,s)}\sech^2 y\in[0,j]$.  Iteration from zero
stays nonnegative, so its unique fixed point $d$ is nonnegative.

To compare this nonnegative mean with the variance, keep $s$ fixed
and define
\[
 H_s(d)=\E_{N(d,s)}[\tanh y-\tanh^2 y].
\]
For $d\ge0$, Gaussian differentiation and
\eqref{eq:gaussian-odd-sign} give
\[
 H_s'(d)=\E_{N(d,s)}[(1-m^2)(1-2m)]\in[-1,1].
\]
Indeed the integrand is bounded below by $-1$, and its expectation is
at most $\E(1-m^2)\le1$ after subtracting
$2\E[m(1-m^2)]\ge0$.  Thus $H_s$ is 1-Lipschitz on $[0,\infty)$.
By \eqref{eq:gaussian-nishimori}, $H_s(s)=0$.  On the other hand,
consistency says $d-s=j(M-q)=jH_s(d)$.  Consequently
\[
 |H_s(d)|\le |d-s|=j|H_s(d)|,
\]
forcing $H_s(d)=0$ and $d=s$.  The consistent Gaussian law must therefore
satisfy \eqref{eq:scalar-fixed-point}.

It remains to solve \eqref{eq:scalar-fixed-point}.  Set
\[
 F(s)=\E_{N(s,s)}\tanh y\quad(s>0),\qquad F(0)=0.
\]
It is continuous at zero by bounded convergence.  For $s>0$,
differentiating both the Gaussian mean and variance gives
\begin{align*}
 F'(s)
 &=\E_{N(s,s)}\left[(\tanh y)'
                       +\tfrac12(\tanh y)''\right]\\
 &=\E_{N(s,s)}[(1-m^2)(1-m)]\in[0,1].
\end{align*}
The lower bound is pointwise; the upper bound follows from
\eqref{eq:gaussian-odd-sign}.  Hence $F$ is nonnegative and
1-Lipschitz on $[0,\infty)$, with $F(s)\le s$.
Consequently, $s\mapsto t+jF(s)$ is a contraction of
$[0,\infty)$ into itself and has exactly one fixed point $s_*(t)$.
The identity \eqref{eq:gaussian-nishimori} then verifies both
consistency equations for $N(s_*(t),s_*(t))$.
If a consistent law has $s=0$, then $0=t+jq$ forces $t=q=0$,
and hence the law is the point mass at zero and $d=0$; this agrees
with the solution just constructed.

Existence and uniqueness are now established, including the
degenerate case.  To obtain the dependence on time, compare the
fixed-point equations at $t,t'$:
\[
 |s_*(t)-s_*(t')|
 \le |t-t'|+j|s_*(t)-s_*(t')|,
\]
which proves the asserted Lipschitz bound.  The size estimates follow
from $0\le F(s)\le s$, which gives $t\le s_*(t)\le t/(1-j)$, and
\eqref{eq:gaussian-nishimori} gives
$q_*(t)=F(s_*(t))=(s_*(t)-t)/j\le t/(1-j)$.
For $t>0$, the variance $s_*(t)$ is positive, so its Gaussian law is
nondegenerate and $q_*(t)>0$.  The equality assertion follows from
Lemma~\ref{lem:scalar-entropy}.
\end{proof}
